%%
%% part1-getting-started/ch02-installation.tex
%%
%% Chapter 2: Installation and Setup — How to install XeLaTeX,
%% fonts, Python dependencies, and MatinBook itself.
%%

\chapter{نصب و راه‌اندازی}
\label{chap:installation}

\section{چرا این موضوع مهم است؟}
\label{sec:installation-why}

قبل از اینکه بتوانید اولین کتاب خود را با
\lr{MatinBook} بسازید، باید چند ابزار را روی
سیستم خود نصب کنید. این ابزارها شامل موتور
حروف‌چینی \lr{XeLaTeX}، فونت‌های فارسی، بسته‌های
\lr{LaTeX}، و چند ابزار جانبی مانند \lr{biber} و
\lr{xindy} هستند.

نصب نادرست یا ناقص، شایع‌ترین دلیل خطاهای اولیه
در \lr{MatinBook} است. بنابراین این فصل را
با دقت بخوانید و هر گام را به‌ترتیب انجام دهید.

\mnote{اگر با \lr{LaTeX} و ابزارهای آن آشنایی دارید، می‌توانید مستقیماً به \cref{sec:installation-matinbook} بروید.}

\section{پیش‌نیازها}
\label{sec:installation-prerequisites}

برای کار با \lr{MatinBook}، به ابزارهای زیر
نیاز دارید:

\begin{enumerate}
    \item \textbf{\lr{TeX Live}:} توزیع اصلی
    \lr{LaTeX} که شامل \lr{XeLaTeX} و اکثر بسته‌ها
    است. نسخه‌ی ۲۰۲۳ یا جدیدتر توصیه می‌شود.

    \item \textbf{\lr{XeLaTeX}:} موتور حروف‌چینی
    که از فونت‌های \lr{OpenType} و حروف‌چینی
    راست‌به‌چپ پشتیبانی می‌کند.

    \item \textbf{\lr{biber}:} ابزار پردازش
    کتاب‌نامه برای \lr{biblatex}.

    \item \textbf{\lr{xindy} و \lr{texindy}:}
    ابزار ساخت نمایه با پشتیبانی از مرتب‌سازی
    فارسی.

    \item \textbf{\lr{Python 3} و \lr{Pygments}:}
    برای رنگ‌آمیزی کد در \lr{minted}.

    \item \textbf{فونت‌های فارسی و لاتین:} شامل
    \lr{XB Niloofar}، \lr{XITS Math}، و
    \lr{DejaVu Sans Mono}.
\end{enumerate}

\mnote{اگر از ویندوز استفاده می‌کنید، نصب \lr{TeX Live} ممکن است چند ساعت طول بکشد. صبور باشید!}

\section{نصب \lr{TeX Live}}
\label{sec:installation-texlive}

\lr{TeX Live} توزیع استاندارد \lr{LaTeX} است
که برای \lr{Linux}، \lr{macOS} و \lr{Windows}
در دسترس است.

\subsection{نصب روی \lr{Linux}}
\label{sec:installation-texlive-linux}

روی \lr{Debian} و \lr{Ubuntu}:

\begin{latin}
\begin{minted}{bash}
sudo apt update
sudo apt install texlive-full
\end{minted}
\end{latin}

\mnote{بسته‌ی \lr{texlive-full} تمام بسته‌های \lr{LaTeX} را نصب می‌کند و حدود ۵ گیگابایت فضا می‌گیرد. اگر فضا محدود دارید، از \lr{texlive-xetex} و بسته‌های موردنیاز دیگر استفاده کنید.}

\subsection{نصب روی \lr{macOS}}
\label{sec:installation-texlive-macos}

روش توصیه‌شده، استفاده از \lr{MacTeX} است:

\begin{latin}
\begin{minted}{bash}
# با Homebrew
brew install --cask mactex
\end{minted}
\end{latin}

\subsection{نصب روی \lr{Windows}}
\label{sec:installation-texlive-windows}

از سایت رسمی \lr{TeX Live}، فایل \lr{install-tl-windows.exe}
را دانلود و اجرا کنید.

\section{نصب فونت‌ها}
\label{sec:installation-fonts}

\lr{MatinBook} به چند فونت نیاز دارد که برخی
از آن‌ها ممکن است روی سیستم شما نصب نباشند.

\subsection{فونت \lr{XB Niloofar}}
\label{sec:installation-font-niloofar}

\lr{XB Niloofar} فونت پیش‌فرض متن فارسی در
\lr{MatinBook} است. اگر در مخزن پروژه هستید:

\begin{latin}
\begin{minted}{bash}
mkdir -p ~/.local/share/fonts
cp fonts/niloofar/*.ttf ~/.local/share/fonts/
fc-cache -fv
\end{minted}
\end{latin}

\mnote{در \lr{macOS}، مسیر \lr{\textasciitilde/Library/Fonts} است. در \lr{Windows}، روی فایل \lr{.ttf} راست‌کلیک کرده و \lr{Install} بزنید.}

برای بررسی نصب:

\begin{latin}
\begin{minted}{bash}
fc-list | grep -i niloofar
\end{minted}
\end{latin}

\subsection{فونت \lr{XITS Math}}
\label{sec:installation-font-xits}

فونت ریاضی \lr{XITS Math} برای نمایش فرمول‌های
ریاضی استفاده می‌شود:

\begin{latin}
\begin{minted}{bash}
# Debian/Ubuntu
sudo apt install fonts-xits

# macOS
brew install --cask font-xits
\end{minted}
\end{latin}

\subsection{فونت \lr{DejaVu Sans Mono}}
\label{sec:installation-font-dejavu}

فونت \lr{DejaVu Sans Mono} برای نمایش کد و
متن فارسی داخل کد استفاده می‌شود:

\begin{latin}
\begin{minted}{bash}
# Debian/Ubuntu
sudo apt install fonts-dejavu
\end{minted}
\end{latin}

\mnote{فونت \lr{DejaVu Sans Mono} روی اکثر توزیع‌های \lr{Linux} به‌طور پیش‌فرض نصب است.}

\subsection{فونت‌های اختیاری}
\label{sec:installation-font-optional}

\lr{MatinBook} از چهار فونت فارسی دیگر نیز
پشتیبانی می‌کند که نصب آن‌ها اختیاری است:

\begin{itemize}
    \item \lr{Vazirmatn}
    \item \lr{Sahel}
    \item \lr{IR Lotus}
    \item \lr{B Nazanin} (تجاری)
\end{itemize}

\section{نصب \lr{Python} و \lr{Pygments}}
\label{sec:installation-python}

برای رنگ‌آمیزی کد با \lr{minted}، به \lr{Python 3}
و \lr{Pygments} نیاز دارید:

\begin{latin}
\begin{minted}{bash}
# Debian/Ubuntu
sudo apt install python3 python3-pip
pip3 install pygments

# macOS
brew install python3
pip3 install pygments
\end{minted}
\end{latin}

برای بررسی نصب:

\begin{latin}
\begin{minted}{bash}
pygmentize -V
\end{minted}
\end{latin}

باید نسخه‌ی \lr{Pygments} را نمایش دهد (مثلاً
\lr{2.17.2}).

\mnote{اگر از \lr{TeX Live 2024} یا جدیدتر استفاده می‌کنید، ممکن است به‌جای \lr{pygments} به \lr{latexminted} هم نیاز داشته باشید. با \lr{pip install latexminted} نصب کنید.}

\section{نصب \lr{MatinBook}}
\label{sec:installation-matinbook}

پس از نصب پیش‌نیازها، نوبت به نصب خود
\lr{MatinBook} می‌رسد. سه روش برای این کار
وجود دارد:

\subsection{روش ۱ — نصب در \lr{texmf} (توصیه‌شده)}
\label{sec:installation-texmf}

این روش، استانداردترین و پایدارترین روش است:

\begin{latin}
\begin{minted}{bash}
mkdir -p ~/texmf/tex/latex/matinbook
cp matinbook.cls ~/texmf/tex/latex/matinbook/
find tex -name "*.sty" -exec cp {} ~/texmf/tex/latex/matinbook/ \;
texhash ~/texmf
\end{minted}
\end{latin}

\mnote{مسیر \lr{\textasciitilde/texmf} در \lr{TeX Live} به‌عنوان \lr{TEXMFHOME} شناخته می‌شود و به‌طور خودکار جستجو می‌شود.}

برای بررسی نصب:

\begin{latin}
\begin{minted}{bash}
kpsewhich matinbook.cls
\end{minted}
\end{latin}

باید مسیر فایل را نمایش دهد.

\subsection{روش ۲ — استفاده از \lr{TEXINPUTS}}
\label{sec:installation-texinputs}

اگر نمی‌خواهید فایل‌ها را در \lr{texmf} کپی کنید،
می‌توانید مسیر پروژه را به \lr{TEXINPUTS} اضافه کنید:

\begin{latin}
\begin{minted}{bash}
export TEXINPUTS="/path/to/matinbook//:"
\end{minted}
\end{latin}

\mnote{دو اسلش \lr{//} یعنی «همه‌ی زیرپوشه‌ها به‌صورت بازگشتی». این روش برای توسعه‌ی پروژه مفید است ولی برای استفاده‌ی نهایی توصیه نمی‌شود.}

\subsection{روش ۳ — نصب در پروژه‌ی محلی}
\label{sec:installation-local}

اگر فقط می‌خواهید در یک پروژه‌ی خاص از
\lr{MatinBook} استفاده کنید، می‌توانید فایل‌ها
را کنار \file{main.tex} قرار دهید. اما این
روش، پایداری کمتری دارد.

\section{تست نصب}
\label{sec:installation-test}

برای اطمینان از نصب درست، یک فایل ساده بسازید:

\begin{latin}
\begin{minted}{latex}
\documentclass{matinbook}

\title{|\pc{تست نصب}|}
\author{|\pc{من}|}
\date{|\pc{مهر ۱۴۰۵}|}

\begin{document}

\maketitle

\chapter{|\pc{فصل اول}|}

|\pc{این یک متن فارسی ساده است.}|

\begin{theorem}
    |\pc{این یک قضیه‌ی آزمایشی است.}|
\end{theorem}

\end{document}
\end{minted}
\end{latin}

\mnote{نام فایل را \lr{test.tex} بگذارید و با \lr{xelatex -shell-escape test.tex} کامپایل کنید.}

سپس با دستور زیر کامپایل کنید:

\begin{latin}
\begin{minted}{bash}
xelatex -shell-escape test.tex
\end{minted}
\end{latin}

اگر یک فایل \lr{test.pdf} ساخته شد و شامل
متن فارسی و قضیه‌ی رنگی بود، نصب شما موفق
بوده است.

\section{عیب‌یابی نصب}
\label{sec:installation-troubleshooting}

در این بخش، چند خطای رایج نصب و راه‌حل آن‌ها
را بررسی می‌کنیم.

\subsection{خطای \lr{File matinbook.cls not found}}
\label{sec:installation-error-cls}

\textbf{علت:} فایل \file{matinbook.cls} در مسیر
جستجوی \lr{LaTeX} نیست.

\textbf{راه‌حل:}
\begin{itemize}
    \item اگر از روش ۱ استفاده می‌کنید،
    \lr{texhash} را اجرا کنید.
    \item اگر از روش ۲ استفاده می‌کنید،
    \lr{TEXINPUTS} را بررسی کنید.
\end{itemize}

\subsection{خطای \lr{Font XB Niloofar not found}}
\label{sec:installation-error-font}

\textbf{علت:} فونت نصب نشده یا \lr{fc-cache}
اجرا نشده است.

\textbf{راه‌حل:}
\begin{latin}
\begin{minted}{bash}
fc-cache -fv
fc-list | grep -i niloofar
\end{minted}
\end{latin}

\subsection{خطای \lr{Package minted Error}}
\label{sec:installation-error-minted}

\textbf{علت:} فراموش کردن \lr{-shell-escape}
یا نصب نبودن \lr{Pygments}.

\textbf{راه‌حل:}
\begin{itemize}
    \item از \lr{xelatex -shell-escape} استفاده کنید.
    \item \lr{Pygments} را نصب کنید:
    \lr{pip3 install pygments}.
\end{itemize}

\subsection{خطای \lr{xindy: command not found}}
\label{sec:installation-error-xindy}

\textbf{علت:} \lr{xindy} نصب نیست.

\textbf{راه‌حل:}
\begin{latin}
\begin{minted}{bash}
sudo apt install xindy
\end{minted}
\end{latin}

\mnote{خطاهای بیشتر در \cref{app:troubleshooting} بررسی می‌شوند.}

\section{تمرین‌ها}
\label{sec:installation-exercises}

\begin{exercise}
    \lr{TeX Live} را روی سیستم خود نصب کنید
    و با دستور \lr{xelatex -v} نسخه‌ی آن را
    بررسی کنید.
    \label{exc:installation-texlive}
\end{exercise}

\begin{exercise}
    سه فونت اصلی \lr{MatinBook} (\lr{XB Niloofar}،
    \lr{XITS Math}، \lr{DejaVu Sans Mono}) را نصب
    کنید و با \lr{fc-list} نصب آن‌ها را تأیید کنید.
    \label{exc:installation-fonts}
\end{exercise}

\begin{exercise}
    یک فایل \lr{test.tex} ساده بسازید و آن را
    با \lr{xelatex -shell-escape} کامپایل کنید.
    اگر خطایی گرفتید، آن را در \cref{sec:installation-troubleshooting}
    جستجو کنید.
    \label{exc:installation-test}
\end{exercise}

\section{خلاصه}
\label{sec:installation-summary}

در این فصل، یاد گرفتیم که:

\begin{itemize}
    \item پیش‌نیازهای \lr{MatinBook} شامل
    \lr{TeX Live}، \lr{biber}، \lr{xindy}،
    \lr{Python} و فونت‌های فارسی و لاتین است.

    \item \lr{TeX Live} روی سه سیستم‌عامل
    \lr{Linux}، \lr{macOS} و \lr{Windows} نصب
    می‌شود.

    \item فونت‌های موردنیاز شامل
    \lr{XB Niloofar}، \lr{XITS Math} و
    \lr{DejaVu Sans Mono} هستند.

    \item \lr{MatinBook} به سه روش نصب می‌شود:
    \lr{texmf} (توصیه‌شده)، \lr{TEXINPUTS} و
    پروژه‌ی محلی.

    \item برای تست نصب، یک فایل ساده بسازید
    و با \lr{xelatex -shell-escape} کامپایل کنید.

    \item خطاهای رایج نصب شامل «فایل پیدا نشد»،
    «فونت پیدا نشد»، «خطای \lr{minted}» و
    «\lr{xindy} پیدا نشد» است.
\end{itemize}

در فصل بعدی (\cref{chap:first-book})، اولین
کتاب خود را با \lr{MatinBook} می‌سازید و
با ساختار پایه‌ی کتاب آشنا می‌شوید.